% Source: source_md/intelligence_part2_appendices_ja.md
\section{創発の VED 内位置づけ}


本補遺は、本論文で導入した「創発の VED 的定義(定義 1.1)」が VED 理論ツリー全体においていかなる位置を占めるかを詳述する。

\bigskip


\subsection{創発概念の扱いの難しさ}


「創発」(emergence)は、複雑系科学・統計力学・科学哲学において多義的に用いられ、論争的な概念である。主要な立場を整理すると:

\begin{longtable}{@{}p{0.30\linewidth}p{0.30\linewidth}p{0.30\linewidth}@{}}
\toprule
立場 & 内容 & 問題 \\
\midrule
\endfirsthead
\toprule
立場 & 内容 & 問題 \\
\midrule
\endhead
弱い創発 & 下位層の記述から原理的に導出可能だが実際上困難 & 計算複雑性の問題に還元される \\
強い創発 & 下位層の記述から原理的に導出不可能 & 二元論的含意・物理的閉鎖性との緊張 \\
全体論的創発 & 全体は部分の和以上 & 「以上」の精密化が困難 \\
機能的創発 & 上位層が下位層とは異なる因果力を持つ & 因果力の帰属の問題 \\
\bottomrule
\end{longtable}


本論文はこれらの論争に直接参入するのではなく、VED の枠組みの内部で「創発」を動力学的・関係的に定義する。

\bigskip


\subsection{VED 的創発の定義(定義 1.1 の再掲)}


\textbf{定義 1.1 (VED 的創発)}

\begin{quote}
ある上位層 $\mathcal{U}$ が、下位層 $\mathcal{L}$ の運動と外部相互拘束 $H_{ij}$ の関数として動力学的に生じるとき、$\mathcal{U}$ は $\mathcal{L}$ から創発するという:

$$
\mathcal{U} = \Phi(\mathcal{L}\text{の運動}, \; H_{ij})
$$
\end{quote}


この定義の三つの含意:

\textbf{含意 1 — 還元主義の否定ではない}:
$\mathcal{U}$ は $\mathcal{L}$ から動力学的に生じる。$\mathcal{L}$ なしに $\mathcal{U}$ はない。これは下位層への依存を認め、完全な自律性を主張しない。

\textbf{含意 2 — 全体論の単純な肯定でもない}:
$\mathcal{U}$ は $\mathcal{L}$ から論理的に導出できない。$H_{ij}$ という外部項が本質的に関与する。下位層の完全な記述だけからは上位層は導けない。これは「強い創発」に近い立場だが、物理的閉鎖性との緊張を回避する——$H_{ij}$ は物理的相互拘束であり、非物理的因果力を仮定しない。

\textbf{含意 3 — 運動的・関係的}:
$\mathcal{U}$ は静的な「性質」や「実体」ではなく、$\mathcal{L}$ の運動が $H_{ij}$ の下で維持される限りにおいて存在する動力学的状態。運動が止まれば $\mathcal{U}$ も消える。

\bigskip


\subsection{VED 内の三例の詳細比較}


\subsubsection{Vol.2: 標準模型構造の創発}


\textbf{文脈}: VED Vol.2 は、因果ログ $C_{ij}$ のネットワークから出発し、閉包エネルギー汎関数 $E(A)$ の最小化を通じて標準模型の構造を導く。

\textbf{創発の機構}:
\begin{itemize}
\item 下位層 $\mathcal{L}$: 因果ログ $C_{ij}$・内部ベクトル $\vec{e}_i$ の運動
\item 外部相互拘束 $H_{ij}$: 閉包条件 $\vec{e}_i + \vec{e}_j + \vec{e}_k \approx 0$(三角形閉包)・閉包エネルギー汎関数 $E(A)$
\item 上位層 $\mathcal{U}$: 粒子(= $E(A)$ の局所極小 = 安定アトラクター)・標準模型のゲージ群構造
\end{itemize}


\textbf{創発の特徴}:
三ノードが形成する最小の自己参照的 enclosure から、粒子が「安定アトラクター」として創発する。これは下位層の運動記述からは論理的に導出されないが、閉包エネルギーという外部拘束の下での動力学的極限として実現される。

\textbf{VED 的創発の観点から}:
粒子の安定性は $E(A)$ の局所極小という動力学的概念であり、静的実体の付与ではない。粒子は「閉包エネルギーの安定アトラクターとして動力学的に維持される状態」であり、VED の運動原理の物理層における現れである。

\subsubsection{Bio-IFGT: 生命の創発}


\textbf{文脈}: Bio-IFGT は、情報場幾何学的な枠組みで化学反応から生命の準閉包が創発する構造を記述する。

\textbf{創発の機構}:
\begin{itemize}
\item 下位層 $\mathcal{L}$: 化学反応・情報流 $J$・状態発展 $\dot{x} = F(x;\theta)$ の運動
\item 外部相互拘束 $H_{ij}$: DNA 由来パラメータ $\theta$・環境との非平衡物質・エネルギー交換
\item 上位層 $\mathcal{U}$: 生命アトラクター $\Omega_{\text{life}}$(準閉包として自己維持する状態)
\end{itemize}


\textbf{創発の特徴}:
化学反応の記述からは生命の自己維持性は論理的に導出されないが、$\theta$(DNA 由来パラメータ)という外部拘束の下での非平衡動力学の極限として創発する。「生命とは準閉包への参入」という VED 的定義が、ここで動力学的に実現される。

\textbf{知性層との接続}:
Bio-IFGT は知性を「生命アトラクターに支えられた第二アトラクター $\Omega_{\text{intel}}$」として予告し、その実装(記憶・言語・個体間共有)を後続研究へ開いた。本論文はこの予告に答える: $L_E$ が第二アトラクターの停止構造であり、$\Phi_{\text{net}}$ が個体間共有の動力学的実体である。

\subsubsection{知性編 Part II: 創発停止層 $L_E$}


\textbf{文脈}: 本論文は推論層 $L_F$ の非閉包性から出発し、ログ層 $L_R$ を経由して創発停止層 $L_E$ が生じる構造を記述する。

\textbf{創発の機構}:
\begin{itemize}
\item 下位層 $\mathcal{L}$: ログ層 $L_R$ の再帰運動
\item 外部相互拘束 $H_{ij}$: 観測者-環境相互拘束($\mu_\theta$ による (B) 組み替えを通じて投影)
\item 上位層 $\mathcal{U}$: 創発停止層 $L_E$(個体内にありながら外部依存的な停止アトラクター)
\end{itemize}


\textbf{創発の特徴}:
$L_R$ の再帰運動は $L_E$ を論理的には含意しない。$\mu_\theta$ による (B) 組み替えが $H_{ij}$ に沿って停止演算子設置領域を投影し、$L_R$ の運動がその設置領域に収束する結果として $L_E$ が動力学的に生じる。

\textbf{Vol.2 との構造的相同}:
Vol.2 において粒子が「$E(A)$ の局所極小 = 安定アトラクター」として創発したように、$L_E$ は「停止汎関数の局所極小 = 停止アトラクター」として創発する。両者は異なる層における同一の創発構造の現れである。

\bigskip


\subsection{VED 的創発と既存創発概念の比較}


\begin{longtable}{@{}p{0.30\linewidth}p{0.30\linewidth}@{}}
\toprule
概念 & VED 的創発との関係 \\
\midrule
\endfirsthead
\toprule
概念 & VED 的創発との関係 \\
\midrule
\endhead
弱い創発(計算複雑性) & 異なる。VED 的創発は計算困難性の問題ではなく、$H_{ij}$ という外部拘束を本質的に要求する構造の問題 \\
強い創発(原理的非還元性) & 部分的に近い。ただし VED 的創発は物理的閉鎖性を維持し、$H_{ij}$ は物理的相互拘束 \\
全体論的創発 & 関係はあるが、VED は「全体」を静的実体として扱わず、「運動の維持状態」として扱う \\
機能的創発 & 部分的に近い。VED 的創発は動力学的機能(停止・生命維持)として記述される \\
\bottomrule
\end{longtable}


VED 的創発は、この中で最も「動力学的・関係的」な立場を取る。上位層は実体ではなく、下位層の運動が外部拘束の下で到達する動力学的状態である。

\bigskip


\subsection{VED 的創発の理論的意義}


VED 的創発概念の理論ツリーにおける意義は三点ある。

\textbf{第一}: 異なる層(物理・生命・知性)を同一の形式 $\mathcal{U} = \Phi(\mathcal{L}, H_{ij})$ で記述することで、領域横断ホモロジーが構造的に担保される。これは「物理から生命から知性」という連続的な理論展開の形式的基盤である。

\textbf{第二}: VED の基底原理「静的範疇ではなく運動」が、創発概念においても貫徹される。創発した上位層は静的性質ではなく、動力学的維持状態として記述される。

\textbf{第三}: 知性層における創発($L_E$)が Vol.2 の粒子創発と同一形式を持つことは、物理から知性まで同一の構造原理が貫通していることの、最も直接的な証拠のひとつである。VED が目指す「単一の原理からの多様な現象の統一的記述」が、創発という概念レベルで実現されている。

\bigskip
